\documentclass[11pt]{article}
\usepackage[margin=0.85in]{geometry}
\usepackage{amsmath,amssymb,amsthm,mathtools}
\usepackage{enumitem}
\usepackage[ruled,vlined,linesnumbered]{algorithm2e}
\usepackage[colorlinks=true,linkcolor=blue]{hyperref}
\usepackage[nameinlink,noabbrev]{cleveref}
\allowdisplaybreaks
\numberwithin{equation}{section}
\newtheorem{assumption}{Assumption}[section]
\newtheorem{definition}[assumption]{Definition}
\newtheorem{theorem}[assumption]{Theorem}
\newtheorem{lemma}[assumption]{Lemma}
\newcommand{\R}{\mathbb R}
\newcommand{\eps}{\epsilon}
\newcommand{\norm}[1]{\left\lVert#1\right\rVert}
\newcommand{\op}{\mathrm{op}}
\newcommand{\Md}{M_{\mathrm d}}
\newcommand{\MP}{M_P}
\DeclareMathOperator*{\argmin}{arg\,min}
\DeclareMathOperator*{\argmax}{arg\,max}
\DeclareMathOperator{\range}{range}
\DeclareMathOperator{\diag}{diag}

\begin{document}

\section{Assumptions, operators, and algorithms}

\begin{assumption}\label{ass:regularity}
Let $f:\R^n\times\R^m\to\R$ be twice continuously differentiable.
There exist $\ell\ge\mu>0$ and $\rho>0$ such that, for all
$\zeta,\zeta'\in\R^{n+m}$,
\begin{equation}\label{eq:regularity}
 \norm{\nabla^2 f(\zeta)}_{\op}\le\ell,
 \qquad
 \norm{\nabla^2 f(\zeta)-\nabla^2 f(\zeta')}_{\op}
 \le\rho\norm{\zeta-\zeta'},
 \qquad
 f_{yy}(\zeta)\preceq-\mu I.
\end{equation}
Moreover,
\[
 P(x):=\max_y f(x,y),\qquad P_*:=\inf_xP(x)>-\infty.
\]
Vector norms are Euclidean and matrix norms are induced operator norms.
The notation $f_x,f_y,f_{xx},f_{xy},f_{yx},f_{yy}$ denotes partial
gradients and Hessian blocks.
\end{assumption}

\begin{definition}\label{def:operators}
For $x\in\R^n$ and $y\in\R^m$, define
\begin{align}
 y^*(x)&:=\argmax_y f(x,y),
 &\kappa&:=\ell/\mu,
 &\Md&:=\rho\kappa^{3/2},
 &\MP&:=4\sqrt2\,\rho\kappa^3,\label{eq:derived}\\
 N_x(y)&:=y-f_{yy}(x,y)^{-1}f_y(x,y),\label{eq:newton-map}\\
 \widehat P(x,y)&:=f(x,y)-\tfrac12 f_y(x,y)^T
                            f_{yy}(x,y)^{-1}f_y(x,y),\label{eq:corrected-value}\\
 \widehat g(x,y)&:=f_x(x,y)-f_{xy}(x,y)f_{yy}(x,y)^{-1}f_y(x,y),\label{eq:corrected-gradient}\\
 \mathcal H(x,y)=\widehat H(x,y)
 &:=f_{xx}(x,y)-f_{xy}(x,y)f_{yy}(x,y)^{-1}f_{yx}(x,y).\label{eq:corrected-hessian}
\end{align}
The cached data at $(x,y)$ are
\[
 \Xi(x,y):=(f(x,y),\nabla f(x,y),\nabla^2f(x,y)),
 \qquad
 \mathcal C(x,y):=(\widehat P(x,y),\widehat g(x,y),\widehat H(x,y);\Xi(x,y)).
\]
An evaluation of each of $f$, $\nabla f$, and $\nabla^2f$ counts as
one query to the corresponding oracle. Cached responses may be reused.
All oracle responses and linear algebra operations are exact.
Trust-region subproblems are solved globally; their arithmetic cost
is not included in the query count.
\end{definition}

\begin{algorithm}[H]
\DontPrintSemicolon
\caption{Persistent Armijo Newton initialization}
\label{alg:initialization}
\KwIn{$f,x_0,y_{\mathrm{init}},\eps>0$}
$\phi(y)\leftarrow-f(x_0,y)$;
 $(y,\alpha)\leftarrow(y_{\mathrm{init}},1)$;
 query and cache $\Xi(x_0,y)$\;
\While{$\norm{\nabla\phi(y)}>\eps^2$}{
 $d\leftarrow-[\nabla^2\phi(y)]^{-1}\nabla\phi(y)$\;
 \While{$\phi(y+\alpha d)>\phi(y)+\frac\alpha2\nabla\phi(y)^Td$}{
  $\alpha\leftarrow\alpha/2$\;
 }
 $y\leftarrow y+\alpha d$;
 query and cache $\Xi(x_0,y)$, reusing the accepted function value\;
}
\Return{$(y,\mathcal C(x_0,y))$}
\end{algorithm}

\begin{algorithm}[H]
\DontPrintSemicolon
\caption{Adaptive fixed-radius trust region with dual Newton tracking}
\label{alg:main}
\KwIn{$f,x_0,y_{\mathrm{init}},\ 0<\eps<1$}
$c\leftarrow2^{-10}$;
 $(x,\sigma)\leftarrow(x_0,1/\log(e/\eps))$\;
$(y,W)\leftarrow\textnormal{Algorithm~\ref{alg:initialization}}
 (f,x_0,y_{\mathrm{init}},\eps)$\;
\While{true}{
 $R\leftarrow c\sqrt\eps/\sigma$;
 $a\leftarrow\sigma\sqrt\eps$\;
 $s\in\argmin_{\norm{u}\le R}
    \{\widehat g^Tu+\tfrac12u^T(\widehat H+aI)u\}$,
 where $(\widehat P,\widehat g,\widehat H)$ are the components of $W$\;
 $\widehat D\leftarrow-\widehat g^Ts-\tfrac12s^T\widehat Hs$\;
 $z\leftarrow x+s$;
 $y_z\leftarrow N_z(y)$;
 $W_z\leftarrow\mathcal C(z,y_z)$\;
 \eIf{$\widehat D>aR^2/4$}{
  \If{$\widehat P-\widehat P_z\ge\widehat D/2$}{
   $(x,y,W)\leftarrow(z,y_z,W_z)$;
   \textbf{continue}\;
  }
 }{
  $\theta\leftarrow\tfrac12\norm{f_{yy}(z,y_z)}_{\op}$;
  $\widehat H_z=\sum_{i=1}^n\lambda_iv_iv_i^T$,
  $\{v_i\}_{i=1}^n$ orthonormal\;
  $p\leftarrow-\sum_{\lambda_i>\theta}
       \lambda_i^{-1}(v_i^T\widehat g_z)v_i$\;
  \If{$\norm{p}\le R$}{
   $w\leftarrow z+p$;
   $y_w\leftarrow N_w(y_z)$;
   $W_w\leftarrow\mathcal C(w,y_w)$\;
   \If{$\norm{\widehat g_w}\le\eps/2$ and
        $\lambda_{\min}(\widehat H_w)\ge-9a/8$}{
    \Return{$w$}
   }
  }
 }
 $\sigma\leftarrow2\sigma$;
 $y\leftarrow N_x(y)$;
 $W\leftarrow\mathcal C(x,y)$\;
}
\end{algorithm}

\section{Convergence and query complexity}

\begin{theorem}\label{thm:main}
Suppose \Cref{ass:regularity} holds. Let
$x_0\in\R^n$, $y_{\mathrm{init}}\in\R^m$, and define
\[
 \Delta_P:=P(x_0)-P_*,\qquad
 \Delta_y:=P(x_0)-f(x_0,y_{\mathrm{init}}).
\]
Assume that $\eps$ satisfies
\begin{subequations}\label{eq:onset}
\begin{align}
 0<\eps&\le e^{-5},
 &\log(e/\eps)&\ge\frac1{\sqrt{\rho\kappa^{3/2}}},\label{eq:onset-scale}\\
 \frac{\eps}{\log^2(e/\eps)}
 &\le\frac{\rho\kappa^2}{2^{19}},\label{eq:onset-init}\\
 \eps\log^2(e/\eps)
 &\le\frac{2^8\mu^2}{\rho^2\kappa^2},\label{eq:onset-radius}\\
 \eps\log^6(e/\eps)
 &\le\frac{2^{24}\mu^2}
              {(1+\kappa)^4\rho^4\kappa^4},\label{eq:onset-hessian}\\
 \eps&\le\min\left\{
       \frac{\mu^2}{64\rho\kappa^{5/2}},
       \frac{8\mu^2}{\rho\kappa^{11/2}}
                         \right\}.\label{eq:onset-polish}
\end{align}
\end{subequations}
Then \Cref{alg:main} terminates at a point $x_{\mathrm{out}}$ satisfying
\begin{equation}\label{eq:stationarity}
 \norm{\nabla P(x_{\mathrm{out}})}<\eps,
 \qquad
 \lambda_{\min}(\nabla^2P(x_{\mathrm{out}}))
 \ge-\frac{145}{64}\sqrt\rho\,\kappa^{3/4}\sqrt\eps
 \ge-\sqrt{\MP\eps}.
\end{equation}
Let $K,J,Q$ be the numbers of accepted trials, failed trials, and
solved trust-region subproblems, respectively. With $c=2^{-10}$,
\begin{align}
 K&\le\frac{16\Delta_P\sqrt\rho\,\kappa^{3/4}}
                 {c^2\eps^{3/2}}+1,\label{eq:K-bound}\\
 J&\le\left\lceil\log_2\left(
             \sqrt\rho\,\kappa^{3/4}\log(e/\eps)\right)\right\rceil,
 \qquad Q=K+J+1.\label{eq:J-bound}
\end{align}
There is a universal numerical constant $C>0$ such that the total
numbers of original-payoff oracle queries satisfy
\begin{equation}\label{eq:query-bound}
\begin{split}
 N_f+N_{\nabla f}+N_{\nabla^2 f}\le C\Bigg[
 &1+\Delta_P\sqrt\rho\,\kappa^{3/4}\eps^{-3/2}
   +\kappa^2\log\left(1+\frac{\ell\Delta_y}{\eps^4}\right)\\
 &+\log\left(1+
          \sqrt\rho\,\kappa^{3/4}\log(e/\eps)\right)\Bigg].
\end{split}
\end{equation}
For every fixed $\ell\ge\mu>0$ and $\rho>0$, the set of accuracies
in \eqref{eq:onset} is a nonempty interval $(0,\eps_0]$.
Its endpoint depends only on $(\ell,\mu,\rho)$, not on $f$, the
dimensions, or the initialization. Neither algorithm takes these
parameters or $\eps_0$ as input. The conclusions
\eqref{eq:stationarity}--\eqref{eq:query-bound} are asserted for
\eqref{eq:onset} only.
\end{theorem}

\section{Value-function geometry}

\begin{lemma}\label{lem:geometry}
Under \Cref{ass:regularity}, $y^*$ is well defined and continuously
differentiable, $P\in C^2$, and
\begin{align}
 Dy^*(x)&=-f_{yy}(x,y^*(x))^{-1}f_{yx}(x,y^*(x)),
 &\norm{Dy^*(x)}_{\op}&\le\kappa,\label{eq:ystar}\\
 \nabla P(x)&=f_x(x,y^*(x)),
 &\nabla^2P(x)&=\mathcal H(x,y^*(x)).\label{eq:P-derivatives}
\end{align}
For every $\zeta=(x,y)$, put
\[
 T(\zeta)=\begin{pmatrix}I\\-f_{yy}(\zeta)^{-1}f_{yx}(\zeta)\end{pmatrix}.
\]
Then
\begin{align}
 -\ell I\preceq\mathcal H(\zeta)&\preceq\frac{\ell^2}{\mu}I,
 \label{eq:schur-spectral}\\
 T(\zeta)^TT(\zeta)&\preceq
       \frac{\mathcal H(\zeta)+(\ell+\mu)I}{\mu},
 \label{eq:relative-geometry}\\
 \norm{\mathcal H(\zeta)-\mathcal H(\zeta')}_{\op}
 &\le(1+\kappa)^2\rho\norm{\zeta-\zeta'},
 \label{eq:schur-lipschitz}\\
 \norm{\nabla^2P(x)-\nabla^2P(x')}_{\op}
 &\le\MP\norm{x-x'}.
 \label{eq:P-lipschitz}
\end{align}
If $\delta=\norm{\nabla^2f(\zeta_1)-\nabla^2f(\zeta_0)}_{\op}$,
then
\begin{equation}\label{eq:schur-comparison}
 \mathcal H(\zeta_1)\succeq
 (1-\delta/\mu)\mathcal H(\zeta_0)
 -\delta(1+\kappa)I.
\end{equation}
\end{lemma}

\begin{proof}
For fixed $x$, $-f(x,\cdot)$ is strongly convex and coercive, so its
minimizer is unique. The implicit function theorem applied to
$f_y(x,y^*(x))=0$ proves \eqref{eq:ystar}. Differentiating
$P(x)=f(x,y^*(x))$ proves \eqref{eq:P-derivatives}.
In particular, $y^*$ is $\kappa$-Lipschitz.

At any $\zeta$, write
\[
 \nabla^2f(\zeta)=
 \begin{pmatrix}A&B\\B^T&-C\end{pmatrix},
 \qquad C\succeq\mu I,
 \qquad \mathcal H=A+BC^{-1}B^T.
\]
The lower bound in \eqref{eq:schur-spectral} follows from
$A\succeq-\ell I$. For its upper bound, fix a unit vector $u$ and
let $v=C^{-1}B^Tu$. If $v=0$, then
$u^T\mathcal H u=u^TAu\le\ell\le\ell^2/\mu$.
If $v\ne0$, set
\[
 e=v/\norm v,\qquad
 \alpha=u^TAu,\quad \beta=u^TBe,\quad \gamma=e^TCe.
\]
The compression to $(u,0),(0,e)$ satisfies
\[
 -\ell I\preceq
 \begin{pmatrix}\alpha&\beta\\\beta&-\gamma\end{pmatrix}
 \preceq\ell I.
\]
Consequently,
\[
 \beta^2\le\min\{(\ell-\alpha)(\ell+\gamma),
                         (\ell+\alpha)(\ell-\gamma)\}.
\]
Since $\beta=\gamma\norm v$,
\[
 \gamma u^T\mathcal H u
 =\gamma\alpha+\beta^2
 \le\ell^2-\ell|\gamma-\alpha|\le\ell^2.
\]
Now $\gamma\ge\mu$ proves the upper bound. Moreover,
\[
 T^TT=I+BC^{-2}B^T
 \preceq I+\mu^{-1}BC^{-1}B^T
 \preceq\mu^{-1}\bigl(\mathcal H+(\ell+\mu)I\bigr),
\]
which proves \eqref{eq:relative-geometry}.

For two points, write $D=f_{yy}(\zeta)$, $D'=f_{yy}(\zeta')$,
$B=f_{xy}(\zeta)$, and $B'=f_{xy}(\zeta')$.
The identity
\[
\begin{split}
 BD^{-1}B^T-B'D'^{-1}B'^T
 ={}&(B-B')D^{-1}B^T
     +B'(D^{-1}-D'^{-1})B^T\\
    &+B'D'^{-1}(B-B')^T
\end{split}
\]
and $D^{-1}-D'^{-1}=D^{-1}(D'-D)D'^{-1}$ give
\eqref{eq:schur-lipschitz}, with coefficient
$\rho(1+2\ell/\mu+\ell^2/\mu^2)$.
Applying it to $(x,y^*(x))$ and $(x',y^*(x'))$ gives
\[
 \operatorname{Lip}(\nabla^2P)
 \le(1+\kappa)^2\sqrt{1+\kappa^2}\rho
 \le4\sqrt2\,\rho\kappa^3.
\]

Finally, for every $u$,
\[
 u^T\mathcal H(\zeta_1)u
 =\max_v\begin{pmatrix}u\\v\end{pmatrix}^{\!T}
              \nabla^2f(\zeta_1)
              \begin{pmatrix}u\\v\end{pmatrix}
 \ge u^T\mathcal H(\zeta_0)u-\delta\norm{T(\zeta_0)u}^2.
\]
Using \eqref{eq:relative-geometry} proves \eqref{eq:schur-comparison}.
\end{proof}

\begin{lemma}[Directional Taylor bounds]\label{lem:directional}
Suppose \Cref{ass:regularity} holds. Let $x,s\in\R^n$ and
$r=\norm s\le\mu/(64\rho\kappa)$. For $t\in[0,1]$, define
\[
 \zeta_t=(x+ts,y^*(x+ts)),\quad
 F_t=\nabla^2f(\zeta_t),\quad
 H_t=\nabla^2P(x+ts),\quad
 T_t=T(\zeta_t),\quad S_t=H_t+(\ell+\mu)I,
\]
\[
 b(t)=\mu^{-1}s^TS_ts
      =(1+\kappa)r^2+\mu^{-1}s^TH_ts.
\]
Then, for $t,u\in[0,1]$,
\begin{equation}\label{eq:path-comparisons}
 b(t)\le4b(0),\qquad S_t\preceq4S_u,
 \qquad\norm{F_1-F_0}_{\op}\le2\rho\sqrt{b(0)}.
\end{equation}
Writing
\[
 \mathcal R=P(x+s)-P(x)-\nabla P(x)^Ts-\tfrac12s^TH_0s,
\]
one has
\begin{equation}\label{eq:scalar-remainder}
 |\mathcal R|\le\tfrac43\rho b(0)^{3/2}.
\end{equation}
In particular,
\begin{align}
 s^TH_0s\le\ell r^2
 &\quad\Longrightarrow\quad
 |\mathcal R|\le4\sqrt3\,\rho\kappa^{3/2}r^3,
 \label{eq:low-curvature-remainder}\\
 s^TH_0s>\ell r^2
 &\quad\Longrightarrow\quad
 |\mathcal R|\le\frac{8\rho\kappa r}{\mu}s^TH_0s.
 \label{eq:high-curvature-remainder}
\end{align}
For $r_g=\nabla P(x+s)-\nabla P(x)-H_0s$,
\begin{equation}\label{eq:full-gradient-remainder}
 \norm{r_g}\le2\sqrt2\,\rho\kappa b(0).
\end{equation}
If $\Pi$ is an orthogonal projector such that
$v^TH_1v\le\ell\norm v^2$ for all $v\in\range\Pi$, then
\begin{equation}\label{eq:projected-gradient-remainder}
 \norm{\Pi r_g}\le2\rho\sqrt{12\kappa}\,b(0).
\end{equation}
\end{lemma}

\begin{proof}
The assertions are immediate if $s=0$. Otherwise $S_t\succeq\mu I$
and $b(t)>0$. The paths $F_t,T_t,H_t$ are Lipschitz. Their derivatives
exist almost everywhere, and
\[
 \norm{F_t'}_{\op}\le\rho\norm{T_ts},\qquad
 F_tT_t=\begin{pmatrix}H_t\\0\end{pmatrix},\qquad
 H_t'=T_t^TF_t'T_t.
\]
Indeed, $H_t=T_t^TF_tT_t$, and both terms involving $T_t'$ vanish
because the upper block of $T_t'$ is zero.
By \eqref{eq:relative-geometry},
\[
 |b'(t)|\le\frac\rho\mu\norm{T_ts}^3
            \le\frac\rho\mu b(t)^{3/2}.
\]
By \eqref{eq:schur-spectral},
\[
 b(0)\le(1+\kappa+\kappa^2)r^2\le3\kappa^2r^2,
 \qquad \frac{\rho\sqrt{b(0)}}{2\mu}\le\frac{\sqrt3}{128}.
\]
Integration of the differential inequality gives
\[
 b(t)\le
 \frac{b(0)}{\bigl(1-t\rho\sqrt{b(0)}/(2\mu)\bigr)^2}
 \le4b(0).
\]
For fixed $v\ne0$, the positive function $v^TS_tv$ satisfies
\[
 \left|\frac{d}{dt}\log(v^TS_tv)\right|
 \le\frac\rho\mu\sqrt{b(t)}
 \le\frac{2\rho\sqrt{b(0)}}\mu
 \le\frac{\sqrt3}{32}<\log4.
\]
Integration proves $S_t\preceq4S_u$. Integration of $F_t'$ proves
the remaining bound in \eqref{eq:path-comparisons}.

The integral Taylor identities are
\[
 \mathcal R=\tfrac12\int_0^1(1-t)^2s^TH_t's\,dt,
 \qquad r_g=\int_0^1(1-t)H_t's\,dt.
\]
Since $|s^TH_t's|\le\rho b(t)^{3/2}$,
\eqref{eq:scalar-remainder} follows. In the low-curvature case,
$b(0)\le(1+2\kappa)r^2\le3\kappa r^2$, proving
\eqref{eq:low-curvature-remainder}. In the high-curvature case,
put $q=s^TH_0s>\ell r^2$. Then
\[
 b(0)\le3q/\mu,\qquad
 q\le\mu\kappa^2r^2,
\]
and hence
\[
 |\mathcal R|\le4\sqrt3\,\rho(q/\mu)^{3/2}
 \le4\sqrt3\,\frac{\rho\kappa r}{\mu}q
 \le8\frac{\rho\kappa r}{\mu}q.
\]

For any unit vector $v$,
\[
 |v^Tr_g|\le
 \rho\int_0^1(1-t)\norm{T_tv}\norm{T_ts}^2\,dt.
\]
The bound $\norm{T_t}_{\op}\le\sqrt{1+\kappa^2}\le\sqrt2\kappa$
proves \eqref{eq:full-gradient-remainder}.
If $v\in\range\Pi$, then
\[
 \norm{T_tv}^2\le\mu^{-1}v^TS_tv
 \le4\mu^{-1}v^TS_1v\le12\kappa.
\]
Taking the supremum over unit vectors in $\range\Pi$ proves
\eqref{eq:projected-gradient-remainder}.
\end{proof}

\section{Newton initialization and oracle errors}

\begin{lemma}[Initialization]\label{lem:initialization}
Under \Cref{ass:regularity}, \Cref{alg:initialization} terminates for
every $\eps>0$ and every $(x_0,y_{\mathrm{init}})$.
Put $\Delta_y=P(x_0)-f(x_0,y_{\mathrm{init}})$.
Let $N_{\mathrm{init}}$ count all its function-value, gradient, and
Hessian queries, including the returned cached tuple.
If $T_0$ is its number of Newton updates and $J_0$
its number of failed Armijo tests, then
\begin{equation}\label{eq:initialization-count}
 T_0\le\left\lceil2\kappa^2\log_+
                   \frac{2\ell\Delta_y}{\eps^4}\right\rceil,
 \qquad J_0\le\lceil\log_2\kappa\rceil,
 \qquad N_{\mathrm{init}}\le3T_0+J_0+3,
\end{equation}
where $\log_+(t):=\log\max\{1,t\}$ for $t\ge0$.
Its returned point satisfies
\begin{equation}\label{eq:initialization-distance}
 \norm{y_0-y^*(x_0)}\le\eps^2/\mu.
\end{equation}
Moreover, for a universal constant $C$,
\begin{equation}\label{eq:initialization-clean-count}
 N_{\mathrm{init}}
 \le C\left[1+\kappa^2\log\left(1+
                                  \frac{\ell\Delta_y}{\eps^4}\right)\right].
\end{equation}
\end{lemma}

\begin{proof}
Write $g=\nabla\phi(y)$, $B=\nabla^2\phi(y)$, and $d=-B^{-1}g$.
Since $\mu I\preceq B\preceq\ell I$, the descent lemma gives
\[
 \phi(y+\alpha d)
 \le\phi(y)+\alpha g^Td+\tfrac12\ell\alpha^2\norm d^2
 \le\phi(y)+\tfrac\alpha2g^Td
 \quad\text{if }\alpha\le\mu/\ell.
\]
Thus every accepted step has $\alpha\ge1/(2\kappa)$, and the
persistent step size is halved at most $\lceil\log_2\kappa\rceil$
times over the entire initialization. Strong convexity implies
$\norm g^2\ge2\mu(\phi(y)-\phi_*)$, and hence
\[
 \phi(y)-\phi(y+\alpha d)
 \ge\tfrac\alpha2g^TB^{-1}g
 \ge\frac1{4\kappa\ell}\norm g^2
 \ge\frac1{2\kappa^2}(\phi(y)-\phi_*).
\]
Also, applying the descent lemma at $y-g/\ell$ yields
$\norm g^2\le2\ell(\phi(y)-\phi_*)$.
Consequently a gap at most $\eps^4/(2\ell)$ suffices for termination,
and geometric contraction proves the bound on $T_0$.
Strong convexity gives
$\mu\norm{y-y^*}\le\norm{\nabla\phi(y)}$, proving
\eqref{eq:initialization-distance}.

The initial cached triple costs three queries; each accepted update
costs one function value and two derivative queries; each failed
Armijo test costs one function value. This proves
\eqref{eq:initialization-count}.
Put $t=\ell\Delta_y/\eps^4$. If $T_0=0$, then $J_0=0$.
Otherwise the initial residual exceeds $\eps^2$, so
$t>1/2$. Therefore $\lceil\log_2\kappa\rceil$
is absorbed into a universal multiple of $\kappa^2\log(1+t)$.
The update bound has the same form, proving
\eqref{eq:initialization-clean-count}.
\end{proof}

\begin{lemma}[Corrected oracle and Newton errors]\label{lem:oracle-errors}
Under \Cref{ass:regularity}, put $e_y=\norm{y-y^*(x)}$.
For every $(x,y)$,
\begin{align}
 \norm{N_x(y)-y^*(x)}&\le\frac\rho{2\mu}e_y^2,
 \label{eq:newton-distance}\\
 \norm{\widehat g(x,y)-\nabla P(x)}
 &\le\frac{(1+\kappa)\rho}{2}e_y^2,
 \label{eq:gradient-error}\\
 \norm{\widehat H(x,y)-\nabla^2P(x)}_{\op}
 &\le(1+\kappa)^2\rho e_y.
 \label{eq:hessian-error}
\end{align}
If $e_y\le\mu/\rho$, then
\begin{equation}\label{eq:value-error}
 |\widehat P(x,y)-P(x)|\le\frac{7\rho}{24}e_y^3.
\end{equation}
\end{lemma}

\begin{proof}
Let $\delta=y-y^*(x)$, $D=f_{yy}(x,y)$, $B=f_{xy}(x,y)$,
and $q=f_y(x,y)$. Taylor's formula gives
\[
 q=D\delta+v,\qquad \norm v\le\tfrac\rho2\norm\delta^2,
\]
\[
 f_x(x,y)-f_x(x,y^*(x))=B\delta+r_x,
 \qquad\norm{r_x}\le\tfrac\rho2\norm\delta^2.
\]
Since $N_x(y)-y^*(x)=-D^{-1}v$,
\eqref{eq:newton-distance} follows. Moreover,
\[
 \widehat g(x,y)-\nabla P(x)=r_x-BD^{-1}v,
\]
and $\norm{BD^{-1}}_{\op}\le\kappa$ proves
\eqref{eq:gradient-error}. Equation \eqref{eq:hessian-error}
is \eqref{eq:schur-lipschitz} at $(x,y)$ and $(x,y^*(x))$.

Taylor expansion of the value at $y$ gives
\[
 P(x)=f(x,y)-q^T\delta+\tfrac12\delta^TD\delta+\mathcal E,
 \qquad |\mathcal E|\le\tfrac\rho6\norm\delta^3.
\]
Substitution of $q=D\delta+v$ yields
\[
 P(x)-\widehat P(x,y)=\mathcal E+\tfrac12v^TD^{-1}v.
\]
Thus
\[
 |P(x)-\widehat P(x,y)|
 \le\frac\rho6e_y^3+\frac{\rho^2}{8\mu}e_y^4
 \le\frac{7\rho}{24}e_y^3
 \quad\text{if }e_y\le\mu/\rho.
\]
\end{proof}

\begin{lemma}[Accuracy at every generated scale]\label{lem:tracking}
Under the hypotheses of \Cref{thm:main}, every stored, trial,
refreshed, and evaluated corrected tuple in \Cref{alg:main},
at its associated scale $\sigma\ge\sigma_0:=1/\log(e/\eps)$,
satisfies
\begin{equation}\label{eq:tracking}
 \norm{y-y^*(x)}\le C_eR^2\le\kappa R,
 \qquad R=c\sqrt\eps/\sigma,
 \qquad C_e:=\frac{2\rho\kappa^2}{\mu}.
\end{equation}
Define
\begin{equation}\label{eq:error-coefficients}
 C_g:=\frac{2(1+\kappa)\rho^3\kappa^4}{\mu^2},
 \quad C_H:=\frac{2(1+\kappa)^2\rho^2\kappa^2}{\mu},
 \quad C_F:=\frac{7\rho^4\kappa^6}{3\mu^3}.
\end{equation}
For each such tuple, define
\[
 \delta_g:=\norm{\widehat g-\nabla P(x)},\qquad
 \delta_H:=\norm{\widehat H-\nabla^2P(x)}_{\op},\qquad
 \delta_F:=|\widehat P-P(x)|.
\]
Then
\begin{equation}\label{eq:strong-errors}
 \delta_g\le C_gR^4,\qquad
 \delta_H\le C_HR^2,\qquad
 \delta_F\le C_FR^6,
\end{equation}
and, with $a=\sigma\sqrt\eps$,
\begin{equation}\label{eq:uniform-errors}
 \delta_g\le\frac{aR}{128},\qquad
 \delta_H\le\frac a{128},\qquad
 \delta_F\le\frac{aR^2}{1024}.
\end{equation}
These assertions hold for every finite execution prefix without
assuming an upper bound on $\sigma$.
\end{lemma}

\begin{proof}
Write $R_0=c\sqrt\eps\log(e/\eps)$.
Conditions \eqref{eq:onset-init} and \eqref{eq:onset-radius} give
\[
 \frac{\eps^2}{\mu}\le C_eR_0^2,
 \qquad R_0\le\frac{\mu}{64\rho\kappa}.
\]
For every $R\le R_0$, $C_eR\le\kappa/32\le\kappa$.
Thus initialization satisfies \eqref{eq:tracking}.
If a primal move has length at most $R$, its warm start has error at
most $\kappa R+\kappa R=2\kappa R$, so one application of
\eqref{eq:newton-distance} gives error at most $C_eR^2$.
This applies to trial moves and, after the explicit norm test, to
the correction $p$.
On a refresh the new radius $R$ equals half the old radius; the old
error is at most $2\kappa R$, and the same argument applies.
The induction does not use an upper scale bound.

In fact $C_eR^2\le\mu/(2048\rho)$, so
\Cref{lem:oracle-errors} applies and gives
\eqref{eq:strong-errors}. By \eqref{eq:onset-hessian},
\[
 \frac{C_HR^2}{a}
 =C_Hc^2\frac{\sqrt\eps}{\sigma^3}
 \le C_Hc^2\sqrt\eps\log^3(e/\eps)\le\frac1{128}.
\]
Furthermore,
\[
 \frac{C_g}{C_H}R
 =\frac{\rho\kappa^2R}{\mu(1+\kappa)}\le\frac1{64},
 \qquad
 \frac{C_F}{C_g}R
 =\frac{7\rho\kappa^2R}{6\mu(1+\kappa)}\le\frac7{384}.
\]
Consequently $C_gR^4\le aR/8192$ and
$C_FR^6\le7aR^2/3145728$, which imply
\eqref{eq:uniform-errors}.
\end{proof}

\section{Acceptance and terminal correction}

\begin{lemma}[Trust-region identities]\label{lem:TR}
Let $a,R>0$, $\widehat H=\widehat H^T$, and let $s$ be a global
minimizer of $\widehat g^Tu+\tfrac12u^T(\widehat H+aI)u$ over
$\norm u\le R$. There exists $\lambda\ge0$ such that
\begin{equation}\label{eq:TR-KKT}
 (\widehat H+(a+\lambda)I)s=-\widehat g,
 \quad\widehat H+(a+\lambda)I\succeq0,
 \quad\lambda(\norm s-R)=0.
\end{equation}
For $\widehat D=-\widehat g^Ts-\tfrac12s^T\widehat Hs$,
\begin{equation}\label{eq:TR-D}
 \widehat D=\tfrac12s^T\widehat Hs+(a+\lambda)\norm s^2
 \ge\tfrac12(a+\lambda)\norm s^2.
\end{equation}
If $\norm{H-\widehat H}_{\op}\le a/128$, then
$s^THs\le2\widehat D$.
If additionally $\widehat D\le aR^2/4$, then
\begin{equation}\label{eq:TR-small}
 \norm s<R,\quad\lambda=0,\quad
 \widehat H\succeq-aI,\quad
 H\succeq-\tfrac{129}{128}aI,\quad
 s^THs\le\tfrac12aR^2.
\end{equation}
\end{lemma}

\begin{proof}
Put $B=\widehat H+aI$. Diagonalize $B=U\diag(\beta_i)U^T$ and
write $g_i=(U^T\widehat g)_i$.
For $\lambda>\max\{0,-\min_i\beta_i\}$, the vector with coordinates
$-g_i/(\beta_i+\lambda)$ has norm continuous and nonincreasing in
$\lambda$, tending to zero at infinity.
If its norm reaches $R$, the resulting vector and multiplier satisfy
\eqref{eq:TR-KKT}. If the limiting norm at the left endpoint is at
most $R$, all singular coordinates have $g_i=0$. At a positive
left endpoint add a nullspace component to make the norm equal to
$R$; at endpoint zero use the limiting vector. This again gives a
feasible pair $(s_*,\lambda_*)$ satisfying \eqref{eq:TR-KKT}.
For any feasible $u$,
\[
 q(u)-q(s_*)
 =\tfrac12(u-s_*)^T(B+\lambda_*I)(u-s_*)
  +\tfrac{\lambda_*}{2}(\norm{s_*}^2-\norm u^2)\ge0.
\]
Thus $s_*$ is a global minimizer. For any other global minimizer,
equality forces the same stationarity and complementarity relations
with $\lambda_*$. Hence \eqref{eq:TR-KKT} holds for the chosen $s$.
Substitution gives \eqref{eq:TR-D}.
Also,
\[
 s^THs\le2\widehat D-2(a+\lambda)\norm s^2
                         +\frac a{128}\norm s^2
 \le2\widehat D.
\]
A boundary solution has $\widehat D\ge aR^2/2$; the small-decrease
condition excludes it. The remaining assertions follow from
\eqref{eq:TR-KKT} and the Hessian error bound.
\end{proof}

\begin{lemma}[Large-decrease acceptance]\label{lem:acceptance}
Under the hypotheses of \Cref{thm:main}, any trial with
$\sigma^2\ge\Md$ and $\widehat D>aR^2/4$ is accepted.
\end{lemma}

\begin{proof}
Let $r=\norm s$, $H=\nabla^2P(x)$, and $q=s^THs$.
If $q\le\ell r^2$, \Cref{lem:directional,lem:TR} give
\[
 |\mathcal R|
 \le4\sqrt3\Md r^3
 \le8\sqrt3\frac{\Md R}{a}\widehat D
 \le8\sqrt3c\widehat D<\frac{\widehat D}{64}.
\]
If $q>\ell r^2$, they give
\[
 |\mathcal R|\le\frac{8\rho\kappa r}{\mu}q
 \le\frac{16\rho\kappa R}{\mu}\widehat D
 \le\frac{\widehat D}{4}.
\]
The working tuples at $x$ and $z=x+s$ satisfy
\eqref{eq:uniform-errors}. Therefore
\[
\begin{aligned}
 \widehat P-\widehat P_z
 &\ge\widehat D-|\mathcal R|
       -\frac{aRr}{128}-\frac{ar^2}{256}
       -\frac{aR^2}{512}\\
 &\ge\left(1-\frac14-\frac1{32}-\frac1{64}-\frac1{128}\right)
           \widehat D
 =\frac{89}{128}\widehat D>\frac12\widehat D.
\end{aligned}
\]
\end{proof}

\begin{lemma}[Small-decrease termination]\label{lem:termination}
Under the hypotheses of \Cref{thm:main}, a trial with
\[
 \sqrt{\Md}\le\sigma\le2\sqrt{\Md},
 \qquad \widehat D\le aR^2/4
\]
passes the norm test on $p$ and both final derivative tests in
\Cref{alg:main}.
\end{lemma}

\begin{proof}
The first bound in \eqref{eq:onset-polish} implies
\begin{equation}\label{eq:a-small}
 a\le2\sqrt{\Md\eps}\le\frac{\mu}{4\sqrt\kappa}\le\ell.
\end{equation}
Let $H_0=\nabla^2P(x)$, $H_1=\nabla^2P(z)$, and use the notation
of \Cref{lem:directional}. By \Cref{lem:TR},
\[
 b(0)\le(1+\kappa)R^2+\frac{aR^2}{2\mu}
 \le\frac52\kappa R^2.
\]
Thus $\delta:=\norm{F_1-F_0}_{\op}\le\sqrt{10}\rho\sqrt\kappa R$.
The radius bound also gives $\delta<\mu$.
Using \eqref{eq:schur-comparison}, \eqref{eq:TR-small}, and
$\sigma^2\ge\Md$,
\begin{equation}\label{eq:endpoint-curvature}
 H_1\succeq
 -\left(1+\frac1{128}\right)aI-\delta(1+\kappa)I
 \succeq-\left(1+\frac1{128}+2\sqrt{10}c\right)aI.
\end{equation}

Let $\Pi_-$ be the spectral projector of $\widehat H_z$ onto
$(-\infty,\theta]$, where
$\theta=\norm{f_{yy}(z,y_z)}_{\op}/2$.
Assumption~\ref{ass:regularity} gives
$\mu/2\le\theta\le\ell/2$. For every $v\in\range\Pi_-$,
\[
 v^TH_1v\le(\theta+a/128)\norm v^2\le\ell\norm v^2.
\]
Equation \eqref{eq:projected-gradient-remainder} therefore applies
to this projector, without any spectral-gap condition.
By the interior stationarity relation in \eqref{eq:TR-KKT},
\[
 \nabla P(x)+H_0s
 =-as+(\nabla P(x)-\widehat g)
       +(H_0-\widehat H)s.
\]
Using the two gradient errors, the initial Hessian error, and the
projected remainder gives
\begin{equation}\label{eq:low-gradient}
\begin{split}
 \norm{\Pi_-\widehat g_z}
 &\le aR+\frac{3aR}{128}
       +10\sqrt3\Md R^2\\
 &\le\left(1+\frac3{128}+10\sqrt3c\right)c\eps
 <2c\eps.
\end{split}
\end{equation}
The full-space bound \eqref{eq:full-gradient-remainder} similarly
gives
\begin{equation}\label{eq:full-endpoint-gradient}
\begin{split}
 \norm{\widehat g_z}
 &\le\left(1+\frac3{128}\right)c\eps
         +5\sqrt2\rho\kappa^2R^2\\
 &\le\left(1+\frac3{128}+5\sqrt2c\right)c\sqrt\kappa\eps
 <2c\sqrt\kappa\eps.
\end{split}
\end{equation}
The correction therefore obeys
\begin{equation}\label{eq:p-bound}
 \norm p\le\theta^{-1}\norm{\widehat g_z}
 \le\frac{4c\sqrt\kappa\eps}{\mu},
 \qquad
 \widehat g_z+\widehat H_zp=\Pi_-\widehat g_z.
\end{equation}
The two conditions in \eqref{eq:onset-polish} are equivalently
\[
 \eps\le\frac{\mu^2}{64\Md\kappa},
 \qquad
 \eps\le\frac{\mu^2\Md}{4096c^2\MP^2\kappa}.
\]
Together with \eqref{eq:p-bound}, they give
\begin{equation}\label{eq:p-small}
 \norm p\le R,
 \qquad \MP\norm p\le a/16.
\end{equation}
Hence the correction is evaluated, and its tuple satisfies
\Cref{lem:tracking}.

Let $w=z+p$. Taylor's formula and \eqref{eq:P-lipschitz} give
\[
\begin{split}
 \widehat g_w
 ={}&\Pi_-\widehat g_z
   +(\nabla P(z)-\widehat g_z)
   +(H_1-\widehat H_z)p\\
   &+r_p+(\widehat g_w-\nabla P(w)),
 \qquad \norm{r_p}\le\tfrac12\MP\norm p^2.
\end{split}
\]
By \eqref{eq:low-gradient}, \eqref{eq:uniform-errors}, and
\eqref{eq:p-small},
\[
 \norm{\widehat g_w}
 \le2c\eps+\frac{3aR}{128}+\frac{aR}{32}
 =\left(2+\frac7{128}\right)c\eps<\eps/2.
\]
For the Hessian, \eqref{eq:endpoint-curvature} gives
\[
 \widehat H_w\succeq
 -\left(1+\frac2{128}+2\sqrt{10}c+\frac1{16}\right)aI
 \succ-\frac98aI.
\]
Both final tests hold.
\end{proof}

\section{Proof of the convergence theorem}

\begin{proof}[Proof of \Cref{thm:main}]
All Newton inverses exist by strong concavity. The initialization
terminates by \Cref{lem:initialization}. Every trust-region subproblem
has a global minimizer by compactness, and $\theta\ge\mu/2>0$
makes each spectral correction well defined. Every other operation
requires finitely many oracle queries.

By \Cref{lem:tracking}, oracle accuracy holds throughout every
finite execution prefix. Put $\sigma_*=\sqrt{\Md}$.
Condition \eqref{eq:onset-scale} gives $\sigma_0\le\sigma_*$.
The first scale at or above $\sigma_*$ lies in
$[\sigma_*,2\sigma_*]$. At such a scale,
\Cref{lem:acceptance,lem:termination} exclude every failure.
It follows, without assuming termination, that
\begin{equation}\label{eq:scale-bound}
 \sigma\le2\sqrt{\Md},\qquad
 J\le\left\lceil\log_2\frac{\sqrt{\Md}}{\sigma_0}\right\rceil
\end{equation}
for every finite prefix.

To bound accepted trials, retain the value-error bound
\begin{equation}\label{eq:E-definition}
 E(\sigma):=C_Fc^6\eps^3\sigma^{-6},
 \qquad E(2\sigma)=E(\sigma)/64.
\end{equation}
Each refresh changes the saved proxy at the same primal point by
at most $E(\sigma)+E(2\sigma)$ in absolute value.
For a finite prefix containing $J$ failures, summing changes of the
saved proxy gives
\begin{align}
 \sum_{\mathrm{accepted}}(\widehat P-\widehat P^+)
 &\le\Delta_P+E(\sigma_0)+E(2^J\sigma_0)
       +\sum_{j=0}^{J-1}\bigl[E(2^j\sigma_0)+E(2^{j+1}\sigma_0)\bigr]
       \notag\\
 &=\Delta_P+2\sum_{j=0}^JE(2^j\sigma_0)
 \le\Delta_P+\frac{128}{63}E(\sigma_0).
 \label{eq:proxy-budget}
\end{align}
This identity uses only saved tuples; rejected trial and correction
tuples do not replace them.
Every accepted trial decreases the saved proxy by more than
\[
 \frac{aR^2}{8}
 =\frac{c^2\eps^{3/2}}{8\sigma}
 \ge\frac{c^2\eps^{3/2}}{16\sqrt{\Md}}.
\]
Thus every finite prefix satisfies
\begin{equation}\label{eq:raw-K-bound}
 K\le\frac{16\sqrt{\Md}}{c^2\eps^{3/2}}
       \left[\Delta_P+\frac{128}{63}E(\sigma_0)\right].
\end{equation}
Applying the coefficient bound $C_gR^4\le aR/128$
from \Cref{lem:tracking} at $\sigma=\sigma_0$ gives
\[
 \eps\log^4(e/\eps)
 \le\frac1{128C_gc^3}
 =\frac{\mu^2}{256c^3(1+\kappa)\rho^3\kappa^4}.
\]
Consequently the error contribution in \eqref{eq:raw-K-bound} is
\begin{align*}
 \frac{2048\sqrt{\Md}}{63c^2\eps^{3/2}}E(\sigma_0)
 &=\frac{2048}{63}\sqrt{\Md}C_Fc^4
             \bigl[\eps\log^4(e/\eps)\bigr]^{3/2}\\
 &\le\frac1{54\sqrt c}
             \frac{\kappa^{3/4}}{(1+\kappa)^{3/2}}
 \le\frac8{27\sqrt2}<1.
\end{align*}
This proves \eqref{eq:K-bound} for every prefix.
Both accepted and failed trials are therefore finite in number.
An infinite execution is impossible, and the algorithm returns.
There is exactly one final terminating trial, so $Q=K+J+1$.

At any return, including one below $\sigma_*$, the all-scale error
bounds and the actual tests imply
\[
 \norm{\nabla P(x_{\mathrm{out}})}
 \le\eps/2+c\eps/128<\eps,
\]
\[
 \lambda_{\min}(\nabla^2P(x_{\mathrm{out}}))
 \ge-\left(\frac98+\frac1{128}\right)\sigma\sqrt\eps
 \ge-\frac{145}{64}\sqrt{\Md\eps}.
\]
Finally,
\[
 \left(\frac{145}{64}\right)^2\Md
 \le4\sqrt2\,\rho\kappa^3=\MP
 \qquad(\kappa\ge1),
\]
which proves \eqref{eq:stationarity}.

Each trial tuple costs at most five queries: a gradient and a
Hessian at the warm start, followed by a full cached triple after
the Newton update. Each evaluated correction costs at most five
queries in the same way. Each refresh costs at most three queries,
because its Newton direction uses the saved derivatives.
There are at most $J+1$ evaluated corrections: each either returns
or causes a failure. Hence
\begin{equation}\label{eq:explicit-query-count}
 N_f+N_{\nabla f}+N_{\nabla^2f}
 \le N_{\mathrm{init}}+5Q+5(J+1)+3J
 =N_{\mathrm{init}}+5K+13J+10.
\end{equation}
Combining \eqref{eq:initialization-clean-count},
\eqref{eq:K-bound}, and \eqref{eq:J-bound} proves
\eqref{eq:query-bound} with a universal constant.

For fixed parameters, all upper-bound left sides in
\eqref{eq:onset} tend to zero as $\eps\downarrow0$, whereas
$\log(e/\eps)$ tends to infinity. On $(0,e^{-5}]$ the functions
$\eps\log^2(e/\eps)$ and $\eps\log^6(e/\eps)$ are nondecreasing,
and $\eps/\log^2(e/\eps)$ is strictly increasing.
Thus the admissible set is a nonempty interval $(0,\eps_0]$
depending only on the common parameters. Inspection of
\Cref{alg:initialization,alg:main} proves the assertion about inputs.
\end{proof}

\end{document}
